\section{Deployment: From Capable Models to Real-Time Systems}
\label{sec:deployment}

The preceding section showed that the evaluation of Speech LLMs must account for multi-channel capabilities and interaction behavior. Deployment introduces a complementary question: whether these capabilities can be delivered continuously under low latency, bounded memory, and realistic hardware constraints. A model that performs well on offline ASR, TTS, or speech understanding benchmarks may still fail as an interactive system if it cannot process streaming input, emit partial output, handle long acoustic contexts, or operate within the resource limits of the target device.

The release of GPT-4o in mid-2024 made this deployment gap broadly visible~\cite{openai2024gpt4o}. It showed that low-latency spoken interaction with multimodal models could approach human conversational timing under the reported conditions: users could speak naturally, receive spoken responses with sub-second latency, and experience an interaction loop that was more fluid than those typically supported by earlier cascaded systems. GPT-4o therefore provides a useful user-experience reference point for the deployment literature: subsequent open, edge, and streaming systems study how parts of this experience can be supported under more transparent or resource-constrained settings. However, reproducing this experience remains difficult outside large proprietary systems. The bottleneck is not only model capability; it is also deployment.

Speech is harder to deploy than text because it is a continuous, multi-channel signal. A deployed Speech LLM must often listen while speaking, preserve linguistic content while tracking paralinguistic cues, produce partial outputs before an utterance is complete, and remain responsive under variable acoustic conditions. These requirements change the serving problem. Batch throughput is no longer sufficient; first-byte latency, token delay, memory growth, turn-taking stability, and edge feasibility become first-order metrics. This section reviews recent work that addresses these constraints through serving systems, deployment-oriented architectures, streaming inference, algorithmic acceleration, and full-duplex omni-modal design.

\subsection{The Deployment Predicament}

Classical speech pipelines separated ASR, dialogue modeling, and TTS. This separation made deployment modular: each component could be optimized and monitored independently. Speech LLMs collapse this boundary. A single model, or a tightly coupled model stack, may now perform recognition, reasoning, response planning, and speech generation. This integration can improve interaction quality, but it also makes deployment failures harder to isolate.

Five recurring deployment failure modes motivate the analysis in this section. The first is latency fragmentation: papers report time to first byte (TTFB), time to first token (TTFT), average token delay, end-to-end latency, or throughput, but these metrics capture different parts of the interaction loop. The second is streaming mismatch: models trained or evaluated on complete utterances are often adapted to streaming use only after training. The third is memory growth: long audio and video streams create key-value cache (KV-cache) and feature-cache pressure that differs from text-only contexts. The fourth is hardware mismatch: theoretical FLOPs do not reliably predict performance on CPUs, mobile NPUs, or mixed edge-cloud deployments. The fifth is interaction incompleteness: many systems optimize recognition or generation latency in isolation, but they do not evaluate simultaneous listening and speaking.

These failure modes show that deployment cannot be summarized by a single latency or throughput number. It depends on how model architecture, decoding policy, memory management, serving infrastructure, and hardware interact under realistic speech workloads.

\subsection{From Failure Modes to Deployment Analysis}

Assessing deployment work requires decomposing each claim into three parts: the operational pressure the system addresses, the system response it applies, and the runtime conditions that produce the result. This structure turns latency, streaming, memory, hardware, and interaction issues into concrete review questions. Does the system reduce first-response delay, incremental token delay, or complete-response time? Does it support partial input during generation? How does memory grow with audio duration? Which device, runtime stack, and serving policy produce the reported result?

The remainder of the section studies five system responses to these pressures: serving-system design, architecture-level deployment design, streaming inference, efficiency optimization, and full-duplex convergence. These responses are not mutually exclusive. A single Speech LLM service may combine a deployment-aware model, a streaming policy, cache and decoding optimizations, a serving stack, and a full-duplex interaction loop.

Table~\ref{tab:deployment-framework} turns this structure into a reporting checklist. The first two columns identify the deployment pressure behind each failure mode. The third column links that pressure to the design patterns most likely to affect it, and the final column specifies the conditions that a deployment card should record. The table is not a leaderboard and does not enumerate deployment papers; it defines the information required to interpret deployment claims in the systems discussed below.

\begin{table*}[t]
\centering
\footnotesize
\setlength{\tabcolsep}{3pt}
\begin{tabularx}{\textwidth}{>{\raggedright\arraybackslash}p{0.17\textwidth} >{\raggedright\arraybackslash}p{0.24\textwidth} >{\raggedright\arraybackslash}p{0.22\textwidth} X}
\toprule
Failure mode & System pressure & Relevant design pattern & Deployment criteria \\
\midrule
Latency fragmentation & First response, incremental output, and complete response measure different behaviors & Serving system; streaming inference & Report TTFB, token delay, end-to-end latency, streaming policy, and interaction-loop latency together. \\
Streaming mismatch & Complete-utterance models may behave differently under partial input & Streaming inference; architecture design & Separate offline and streaming settings; specify input chunk size, output granularity, and emission policy. \\
Memory growth & Long audio and video streams expand KV-cache, feature cache, and multimodal state & Efficiency optimization; serving system & Report input duration, output duration, cache policy, peak memory, and active compression or eviction methods. \\
Hardware mismatch & FLOPs and parameter count do not determine device-level behavior & Architecture design; serving system & Report device class, runtime stack, batch size, measured latency, measured memory, and edge/cloud placement. \\
Interaction incompleteness & Recognition or generation latency alone does not capture simultaneous listening and speaking & Full-duplex convergence; serving system & Evaluate simultaneous input/output, turn-taking, barge-in, interruption recovery latency, and response quality together. \\
\bottomrule
\end{tabularx}
\caption{A deployment analysis framework for Speech LLMs. Each failure mode induces a system pressure, points to relevant design patterns, and implies concrete criteria for deployment-card reporting.}
\label{tab:deployment-framework}
\end{table*}

The corresponding subsections discuss the quantitative signals reported by these systems. We do not aggregate them into a leaderboard because the systems differ in task, hardware, latency definition, and interaction setting.

\subsection{System-Level Serving}

System-level serving determines whether Speech LLM inference can run as temporal-stream execution rather than isolated prompt processing. The central design question is how to coordinate input ingestion, acoustic encoding, language generation, speech synthesis, customization, and cluster scheduling under one serving loop while controlling latency fragmentation, memory growth, hardware mismatch, and interaction incompleteness.

VoxServe (Kamahori et al., 2026) introduces a serving system designed specifically for Speech Language Models~\cite{kamahori2026voxserve}. Its main architectural idea is a model-execution abstraction that separates model architecture from system-level optimization. Modern SpeechLMs vary widely in their internal structures, including encoder-decoder ASR models, autoregressive speech generators, and codec-based synthesizers. Nevertheless, they share computational patterns such as attention over temporal sequences, streaming I/O, and multi-stage processing. VoxServe exploits these shared patterns to support streaming-aware scheduling and an asynchronous inference pipeline, reporting 10--20x higher throughput than existing implementations at similar latency across several modern SpeechLMs.

NIM4-ASR (Xie et al., 2026) studies serving from the perspective of model-system co-design~\cite{xie2026nim4asr}. Instead of asking how to serve an arbitrarily large model efficiently, it asks what model size is sufficient for production-grade quality and how such a model should be served at scale. The resulting system uses a 2.3B-parameter encoder-LLM architecture with a clear division of labor: the encoder handles acoustic processing, while the LLM handles linguistic reasoning. The production system supports real-time streaming inference and million-scale hotword customization through retrieval-augmented generation (RAG) with sub-millisecond retrieval latency.

EPD Disaggregation (Singh et al., 2025) extends the prefill-decode disaggregation paradigm from text LLMs to multimodal serving~\cite{singh2025epd}. Its main observation is that multimodal models introduce a separate bottleneck: encoding multimedia inputs such as video frames and audio spectrograms is expensive and operates on a different timescale from prefilling or token generation. The system separates serving into three independently scheduled phases, Encode, Prefill, and Decode, and adds a caching mechanism for multimedia tokens.

Evaluating Kubernetes Performance for GenAI Inference (Malleni et al., 2026) studies the infrastructure layer and reports substantial improvements in tail TTFT latency under Kubernetes-native routing and scheduling configurations, with the exact percentage depending on the metric and experimental setting~\cite{malleni2026kubernetes}. Although this work is not specific to Speech LLMs, it is relevant because speech interaction is highly sensitive to tail latency.

The shared pattern is that serving decisions determine which part of the speech interaction loop becomes visible. VoxServe emphasizes heterogeneous temporal execution, NIM4-ASR emphasizes the production ASR boundary between acoustic encoding and linguistic reasoning, EPD Disaggregation emphasizes phase separation, and Kubernetes-level work emphasizes cluster behavior. Speech LLM deployment needs serving stacks that expose throughput, TTFT, and interaction-level traces under the same setup, including simultaneous input ingestion, output generation, interruption handling, cache growth, and tail latency under realistic audio durations.

\subsection{Architecture-Level Deployment Design}

Architecture-level work brings deployment constraints into the model design stage. Instead of treating latency, memory, and hardware behavior only as serving issues, these systems make them part of the architecture itself. This premise is especially important for edge and consumer hardware, where memory and latency limits cannot be addressed only through better serving.

LFM2 (Amini et al., 2025) uses hardware-in-the-loop architecture search under explicit edge latency and memory constraints~\cite{amini2025lfm2}. The resulting hybrid backbone combines gated short convolutions with grouped query attention blocks and is optimized for measured performance on target hardware rather than theoretical FLOPs. The model family ranges from 350M to 8.3B parameters, includes dense and MoE variants, supports 32K context length, and is trained on 10--12 trillion tokens. The LFM2-Audio variant supports real-time speech-to-speech interaction, reports competitive quality relative to substantially larger models under its evaluation setting, and achieves 2x faster prefill and decode on CPUs than similarly sized models.

MiniCPM-o 4.5 (Cui et al., 2026) focuses on full-duplex interaction~\cite{cui2026minicpmo}. Many multimodal models still operate in alternating phases: they perceive and then respond. This design prevents the model from incorporating new inputs during generation or acting proactively. MiniCPM-o introduces Omni-Flow, a streaming framework that aligns omni-modal inputs and outputs on a shared temporal axis and transforms turn-based interaction into continuous, time-aligned processing. With 9B parameters and less than 12GB RAM, it is presented as an open-source model that supports simultaneous seeing, listening, and speaking with proactive behavior on consumer edge hardware.

S2M3 (Yoon et al., 2025) takes a modular approach to edge deployment by decomposing multi-modal models into reusable modules for distributed multi-task inference~\cite{yoon2025s2m3}. By allowing the same encoder modules to serve multiple downstream tasks, S2M3 reduces memory use by up to 62\% and improves latency by 56.9\% relative to cloud-based approaches. TinyM2Net-V3 (Rashid \& Mohsenin, 2024) studies a more constrained setting, combining knowledge distillation with low-bit-width quantization for multimodal audio-visual inference under severe memory limits~\cite{rashid2024tinym2net}.

These systems represent four distinct deployment philosophies. LFM2 optimizes directly for measured hardware behavior, MiniCPM-o optimizes for interaction continuity, S2M3 optimizes for modular resource sharing, and TinyM2Net-V3 targets severe memory constraints through compression. Together, they show that deployment-aware architecture design spans trade-offs among latency, memory, modularity, and interaction capability. Speech LLM deployment needs architecture specifications that include device class, measured latency, memory footprint, supported context length, input audio duration, and output audio duration.

\subsection{Streaming and Real-Time Inference}

Streaming work makes output timing part of the inference objective. A streaming model must decide not only what to output, but also when to output it. Component-level latency and conversational-loop latency are related but not identical: a correct response that arrives too late may still fail the interaction.

Ultra-Low Latency End-to-End Streaming Speech Synthesis (Su et al., 2026) reports 48.99ms average TTFB and 10.6x acceleration over conventional cascaded pipelines under its evaluation setting~\cite{su2026ultralowtts}. The key idea is to remove the vocoder bottleneck by directly modeling the compressed discrete latent space of the Mimi neural audio codec. A modified FastSpeech 2 backbone uses progressive depth-wise sequential decoding across 32 residual vector quantization layers to generate speech in non-autoregressive blocks. Experiments in English and Malay provide cross-language evidence within the scope of the evaluated languages.

SpeakStream (Bai et al., 2025) addresses the LLM-then-TTS cascade used in many conversational AI systems~\cite{bai2025speakstream}. Traditional TTS models are trained on complete utterances and can introduce large delays when paired with streaming LLM text. SpeakStream uses a decoder-only architecture trained with next-step prediction on interleaved text-speech data, enabling it to generate audio incrementally while receiving streaming text.

Streaming ASR with Decoder-Only LLMs (Wan et al., 2026) studies when an LLM-based ASR system should emit tokens relative to the input speech~\cite{wan2026streamingasr}. A read/write policy network with monotonic chunkwise attention reduces average token generation delay by 62.5\% under the reported setup. Mini-Omni (Xie \& Wu, 2024) was an early end-to-end system for real-time speech interaction without an external TTS system~\cite{xie2024miniomni}. Its text-instructed speech generation method, batch-parallel inference strategy, and ``Any Model Can Talk'' training method showed that a single model could handle the complete speech interaction loop while preserving core language capabilities.

The shared pattern is that different systems optimize different latency boundaries. TTFB, first-token latency, token generation delay, and end-to-end latency answer different questions. A speech synthesis result with strong TTFB and an ASR result with low token delay provide complementary views of latency, but neither metric alone characterizes the full interaction loop. Speech LLM deployment needs latency optimization at the conversational-loop level, where component latency links to streaming policy, input chunk size, output granularity, and downstream user-visible response time.

\subsection{Efficiency Optimization}

Efficiency optimization treats decoding, compression, and cache management as deployment levers that can control memory growth and hardware mismatch without redesigning the entire model or serving stack. Real systems may combine quantization, speculative decoding, cache pruning, batching, and streaming policies, so the important deployment question is how these gains interact when they are composed.

\subsubsection{Compression and Quantization}

Quantizing Whisper-small (Sohler et al., 2026) evaluates post-training quantization across libraries and bit-widths for Whisper and provides practical guidance for constrained hardware~\cite{sohler2025quantizingwhisper}. Edge-ASR (Feng et al., 2025) expands this analysis to eight post-training quantization (PTQ) methods across Whisper and Moonshine, and identifies best practices for low-bit edge deployment~\cite{feng2025edgeasr}. Resource-Efficient Speech Quality Prediction (Nilsson et al., 2024) shows that binary activation maps with 8-bit quantization-aware training (QAT) can reduce inference memory by 25-fold for speech quality assessment~\cite{nilsson2024speechquality}. These results suggest that speech models may tolerate aggressive quantization in some settings, but this tolerance depends on the task, architecture, and evaluation metric.

\subsubsection{Speculative Decoding}

Speculative decoding uses a fast draft model to propose tokens that a larger verifier accepts or rejects in parallel. SpecASR (Wei et al., 2025) reports 3.04--3.79x speedup over autoregressive decoding with no accuracy loss under its evaluation setting, using adaptive draft generation and a sparse token tree~\cite{wei2025specasr}. Accelerating Autoregressive Speech Synthesis (Lin et al., 2025) applies speculative decoding to TTS and reports a 1.4x speedup while maintaining fidelity and naturalness~\cite{lin2025speechspeculative}. The smaller gain relative to ASR is consistent with the higher token entropy of speech synthesis.

Edge-Cloud Collaborative Speech Emotion Captioning (Xue et al., 2026) introduces Uncertainty-Guided Speculative Decoding for privacy-preserving edge-cloud collaboration~\cite{xue2026edgecloudsec}. A lightweight edge model drafts emotion captions locally, and only high-uncertainty token blocks are sent to a cloud verifier. Model-free Speculative Decoding (Ho et al., 2025) removes the draft model and instead uses precomputed n-gram token maps, reporting a 10\% absolute speed improvement on CPU with minimal additional memory~\cite{ho2025modelfreeasr}.

\subsubsection{KV-Cache Optimization}

AccKV (Jiang et al., 2025) addresses KV-cache growth in audio-video LLMs, where both modalities create long temporal sequences~\cite{jiang2025acckv}. The key observation is that attention in higher transformer layers shifts toward the video modality regardless of task type. AccKV uses Layer Adaptive Focusing to identify important tokens by layer and modality, and Cross-Calibration to align low-priority modality caches with high-priority ones before selective eviction.

Across these studies, the shared pattern is modular efficiency: each method improves one pressure point in the deployment stack. Speech LLM deployment needs to treat these methods as composable system policies rather than isolated speedups. For example, evaluations can combine quantization, speculative decoding, and cache pruning under the same streaming input regime. This is the deployment analogue of evaluating individual benchmark capabilities together rather than only in isolation.

\subsection{Deployment Evidence Under Realistic Pressure}

Deployment evidence should specify whether a result comes from a product demonstration, a reproducible component benchmark, or a controlled system measurement. GPT-4o-like products define compelling user-facing targets, while open systems make instrumentation, ablations, and controlled measurement easier. Researchers often study Whisper-derived ASR models, codec-based TTS pipelines, and selected open multimodal models because these systems support reproducible instrumentation. Interaction-first systems such as MiniCPM-o 4.5, OpenOmni, GPT-4o-Audio-style real-time systems, and Gemini Live-style products provide references for interaction quality, although access, instrumentation, and reproducibility differ across systems.

Speech LLM deployment needs to separate three kinds of evidence: product-level demonstrations, reproducible component benchmarks, and controlled system measurements. Product demonstrations define experience targets; component benchmarks isolate mechanisms; controlled system measurements compare deployment behavior. Treating these as complementary forms of evidence makes cross-system comparison clearer without using one form as a substitute for another.

\subsection{Full-Duplex and Omni-Modal Convergence}

Full-duplex omni-modal systems treat perception and generation as concurrent rather than alternating processes. Full-duplex interaction is therefore not merely faster streaming; it changes the interaction contract. A system must process incoming audio or visual streams while generating speech, maintain low latency in both directions, accept interruptions, and coordinate multiple modalities under memory and compute constraints.

OpenOmni (Luo et al., 2025) introduces a two-stage progressive multimodal alignment framework with 7B parameters~\cite{luo2025openomni}. A notable result is cross-modal transfer: training a pre-trained speech model on text-image tasks yields near-zero-shot generalization from vision to speech and outperforms models trained on explicit tri-modal datasets. A lightweight non-autoregressive decoder trained with direct preference optimization produces real-time emotional speech with sub-1-second latency and a 5x speedup over autoregressive methods. Together with MiniCPM-o 4.5, OpenOmni illustrates a broader trend toward omni-modal models in the 7--9B range that can run on consumer hardware with real-time responsiveness.

This convergence changes the deployment target. The field no longer only needs faster ASR, faster TTS, or cheaper serving; it needs an interaction loop in which perception, reasoning, and speech generation occur continuously. Speech LLM deployment should treat full-duplex interaction as a first-class target and measure turn-taking, barge-in, interruption recovery latency, memory growth, and content quality together because these properties determine whether a system remains responsive and coherent during overlapping input and output.

\subsection{Deployment Recommendations}

Speech LLM deployment should design, optimize, and describe systems as real-time interaction loops rather than as independent ASR, LLM, and TTS components. Speculative decoding, quantization, KV-cache pruning, streaming serving, and edge deployment all improve different pressure points, but their value is clearest when they are tied to the deployed interaction regime.

First, deployment should move from isolated optimization to composed optimization. When a system combines quantization, speculative decoding, cache pruning, batching, or streaming policies, the evaluation should treat these choices as one deployment configuration whose combined behavior determines latency, quality, and memory growth.

Second, deployment should make streaming state explicit. Input chunk size, output granularity, cache policy, audio duration, interruption setting, and simultaneous input-output status define the runtime state of a Speech LLM. These details connect component speedups to the actual interaction loop.

Third, deployment should align latency metrics with interaction conditions. TTFB, average token delay, end-to-end latency, and time-to-completion each characterize a different user-visible behavior: first response, incremental output, complete answer delivery, or interruption recovery.

A deployment card summarizes these recommendations as a compact disclosure of the operational conditions under which an evaluation measures a deployment claim. At minimum, it should include model size, hardware, batch size, streaming policy, input duration, output duration, active optimizations, latency metric, quality metric, baseline implementation, and whether the evaluation uses cloud, edge, streaming, or full-duplex settings. This practice makes cross-paper comparison more tractable even when tasks differ.

The broader recommendation is to report both component-level measurements and interaction-loop measurements. A low TTFB becomes more interpretable when the evaluation connects it to conversational responsiveness. Similarly, evaluations should tie ASR token delay to downstream response delay when the target application is an assistant or agent. For full-duplex systems, evaluation should measure turn-taking, interruption handling, latency, and content quality together. This is the deployment counterpart of multi-channel evaluation: the system must be correct, responsive, and stable at the same time.

These deployment choices also determine the safety surface of Speech LLMs. Systems that stream continuously, retain acoustic context, split computation across edge and cloud, or synthesize speech in real time expose new risks related to privacy, impersonation, provenance, and audio-specific attacks. Section~\ref{sec:safety_privacy_governance} turns to these safety, privacy, and governance implications.
